\section{模型评价与结论}
本文围绕三个递进目标，分别把节点顺序、物理空间和流水等待显式化。生命周期收缩在固定操作顺序下不增峰值；连续分配通过完整驻留状态与有代价的恢复机制获得可行方案；固定驻留图重排保持地址先后和搬运事件不变，从而将时间优化转化为受资源约束的拓扑排序。

六图中，问题一只有卷积小图相对收缩后的编号基准降低峰值，降幅为 7.62\%。问题二的搬运量依次为 44766、75546、3620、32852、34944 和 460800。问题三在这些搬运量均不增加的条件下，将总时间降低 1.30\% 至 33.38\%；其中注意力大图从 225902 降至 150499 周期。阶段比较表明，四图的改善来自固定驻留重排，两个矩阵乘图还需要先改变地址分配候选。

在当前六图与驻留段模型下，缓存方案确定后仍存在可利用的流水优化空间。卷积小图距固定驻留下界为 7040 周期，进一步研究可以转向地址关系；卷积大图与该下界相差 24.95\%，适合继续扩展单元顺序搜索。矩阵乘大图的输入搬运工作量为 1628000 周期，可优先研究重复换出对象的使用组织。不同方向分别对应地址链、资源排序与搬运事件三个层次。

过程分析进一步给出这些方向的具体依据。卷积小图的选定峰值由一级缓存 6912 与零级右输入缓存 72 组成，局部缓存贪心却在零级右输入缓存积累了 16384 的逻辑需求；矩阵乘大图的 3600 次换出涉及 255 个不同对象，反复恢复是当前方案的重要开销；注意力小图在固定搬运量条件下，多处向量节点的依赖等待缩短并向后传播，最终完成窗口减少 9765 周期。这些结果把汇总指标落实到实际对象、事件及时间窗口。

方法的主要价值是将三个目标组织为可逐层检验的完整方案：生命周期收缩提供统一基准，连续分配保证物理空间与恢复状态合法，固定驻留调度利用地址链保留的并行机会。全局峰值和搬运量的紧下界、不同图结构上的泛化表现以及真实芯片对应关系，构成后续研究的三个重点。现阶段结论适用于题面六图及本文明确的驻留段复用解释。
